% ---
% id: protocol-compiler/rem:wellposed
% kind: rem
% title: "Hermiticity as a declared condition"
% summary: >-
%   The left-hand sides of both constraint lines are Hermitian, respectively
%   real, because the constraint operators and ω are Hermitian
%   ([[lem:hermitian-constraint-values]]). Without Hermiticity the lines lose
%   their meaning, which is why it is a condition of the declaration, decided
%   exactly.
% prominence: standard
% status: accepted
% uses: [def:feasible, def:declaration, lem:hermitian-constraint-values]
% mentions: [def:validator]
% assumes: []
% refs: []
% provenance: newly-derived
% lean:
% ---

\begin{remark}
\label{rem:wellposed}
By Lemma~\ref{lem:hermitian-constraint-values} both sides of the trusted line
of \eqref{eq:feas-modes} are Hermitian, so that their L\"owner comparison is
defined, and the left-hand side of the untrusted line is real, so that it can
be compared with $c_k\in\mathbb{R}$. Drop Hermiticity and both lines lose
their meaning rather than merely their soundness, which is why it is a
condition of Definition~\ref{def:declaration} and of D(iii), and not only a
check.

\emph{Where each hypothesis of that lemma comes from.} $X=X^{\dagger}$, for
$X=\idon{A}\otimes\hat O_k$, is D(iii)'s Hermiticity clause on $\hat O_k$,
decided by check V3 of the validator, which decides the operators block by
block (Definition~\ref{def:validator}). $\omega=\omega^{\dagger}$ is the
Hermiticity half of the positivity condition of \eqref{eq:feas}, decided by
the feasibility check V7 on the candidate it verifies: a condition on
$\omega$ and not an assumption about it. Both are decided exactly.
\end{remark}
